\chapter{Introduzione}
\label{cap:introduzione}

\section{Scopo del documento}

\progetto{} è un servizio minimo di gestione ordini: riceve un ordine e lo salva in un database. Il servizio viene realizzato in cinque \emph{architetture} (V1--V5), sottoposto a guasti controllati e osservato mentre si ripristina. Per ogni architettura si misurano il tempo di fermo, gli ordini falliti, quelli persi e il tempo necessario a tornare operativi.

Il servizio è volutamente semplice: oggetto del progetto non è il codice, ma le \textbf{scelte architetturali} e il loro effetto in presenza di guasti. Questo documento descrive le architetture, gli ambienti di esecuzione, il modello dei guasti, le decisioni prese e il protocollo di misura, e raccoglie le previsioni formulate prima degli esperimenti.

Il nome viene dal \textit{kintsugi}, l'arte giapponese di riparare la ceramica rotta con l'oro, mettendo in evidenza le crepe invece di nasconderle: ogni guasto misurato mostra dove l'architettura cede e come si ricompone.

\section{Metodo}

Il progetto segue tre regole.

\begin{enumerate}
  \item \textbf{Un fattore alla volta.} Ogni architettura cambia un solo aspetto rispetto a quella da cui deriva, così ogni differenza misurata ha una causa attribuibile.
  \item \textbf{Prima prevedere, poi misurare.} Le previsioni si scrivono e si registrano su Git \emph{prima} degli esperimenti (tag \texttt{m1}). Lo scarto tra previsto e misurato indica quanto il comportamento del sistema è stato compreso.
  \item \textbf{Dichiarare i limiti.} Ciò che gli esperimenti non possono dimostrare è dichiarato esplicitamente (capitolo~\ref{cap:validita}).
\end{enumerate}

La figura~\ref{fig:metodo} riassume il ciclo di lavoro.

\begin{figure}[htbp]
\centering
\begin{tikzpicture}[node distance=0.55cm,
  fase/.style={box, minimum width=2.45cm, minimum height=1.25cm},
  doc/.style={etichetta, text width=2.45cm}]
\node[fase] (p) {\textbf{Prevedere}\\{\scriptsize modelli e parametri}};
\node[fase, right=of p] (c) {\textbf{Configurare}\\{\scriptsize una versione V1--V5}};
\node[fase, right=of c, guastato] (g) {\textbf{Guastare}\\{\scriptsize un componente F1--F7}};
\node[fase, right=of g] (m) {\textbf{Misurare}\\{\scriptsize client, DB, cluster}};
\node[fase, right=of m] (k) {\textbf{Confrontare}\\{\scriptsize previsto e misurato}};
\foreach \a/\b in {p/c, c/g, g/m, m/k} \draw[freccia] (\a) -- (\b);
\node[doc, below=2mm of p] {cap.~\ref{cap:previsioni}\\tag \texttt{m1}};
\node[doc, below=2mm of c] {cap.~\ref{cap:architettura}, \ref{cap:deployment}};
\node[doc, below=2mm of g] {cap.~\ref{cap:guasti}};
\node[doc, below=2mm of m] {cap.~\ref{cap:misura}};
\node[doc, below=2mm of k] {report finale\\(RQ1--RQ6)};
\draw[risposta] (k.north) -- ++(0,0.7) -| node[etichetta, pos=0.25, above] {lo scarto spiega cosa non avevamo capito del sistema} (p.north);
\end{tikzpicture}
\caption{Il metodo del progetto: ogni esperimento confronta una previsione scritta in anticipo con una misura.}
\label{fig:metodo}
\end{figure}

\section{Struttura del documento}

\begin{table}[htbp]
\centering
\caption{Organizzazione dei capitoli.}
\label{tab:struttura}
\begin{tabularx}{\linewidth}{@{}lL@{}}
\toprule
\textbf{Capitolo} & \textbf{Contenuto} \\
\midrule
\ref{cap:obiettivi} & Obiettivo, domande di ricerca, ipotesi e ambito \\
\ref{cap:architettura} & Le cinque architetture, il percorso di un ordine e i punti unici di guasto \\
\ref{cap:deployment} & Gli ambienti di esecuzione: cluster locale e AWS \\
\ref{cap:guasti} & Il modello dei guasti e il modo in cui vengono iniettati \\
\ref{cap:decisioni} & Il registro delle decisioni architetturali (ADR) e i parametri \\
\ref{cap:misura} & Definizioni, struttura di un run, carichi, matrice e statistica \\
\ref{cap:previsioni} & Modello di previsione e previsioni \\
\ref{cap:validita} & Minacce alla validità e contromisure \\
Appendici & Pianificazione e correzioni apportate al piano \\
\bottomrule
\end{tabularx}
\end{table}
